\chapter{Outils numériques : schéma d'Euler explicite}

\section{Objectif de la résolution numérique}
De nombreux modèles de sciences de l'ingénieur conduisent à une équation différentielle dont la solution exacte n'est pas directement exploitable. Une méthode numérique construit alors une suite de valeurs approchées de la solution à des instants discrétisés.

On considère le problème de Cauchy
\[
\dot y(t)=f\bigl(t,y(t)\bigr),\qquad y(t_0)=y_0.
\]
On choisit un pas de temps constant $h>0$ et les instants
\[
t_k=t_0+kh.
\]
La valeur calculée $y_k$ constitue une approximation de $y(t_k)$.

\begin{center}
\TLEulerGeometrique
\end{center}

\section{Construction du schéma d'Euler explicite}
Au voisinage de $t_k$, la dérivée peut être approchée par
\[
\dot y(t_k)\simeq \frac{y(t_{k+1})-y(t_k)}{h}.
\]
Comme $\dot y(t_k)=f(t_k,y(t_k))$, on obtient la relation de récurrence
\[
\boxed{y_{k+1}=y_k+h f(t_k,y_k)}.
\]
Cette relation est dite \emph{explicite} car la nouvelle valeur $y_{k+1}$ est calculée uniquement à partir de valeurs déjà connues.

\begin{center}
\TLAlgorithmeEuler
\end{center}

\section{Exemple linéaire du premier ordre}
Les sources d'exercices proposent l'équation
\[
\dot y(t)+\alpha y(t)=\beta,\qquad y(0)=\gamma.
\]
Elle s'écrit
\[
\dot y(t)=\beta-\alpha y(t).
\]
Le schéma d'Euler explicite donne alors
\[
\boxed{y_{k+1}=y_k+h\left(\beta-\alpha y_k\right)}
\]
soit
\[
y_{k+1}=(1-\alpha h)y_k+\beta h.
\]
Le facteur $\alpha h$ montre que le choix du pas influence directement la qualité et la stabilité du calcul.

\section{Équation non linéaire du premier ordre}
Pour
\[
\dot y(t)=-t\,y^2(t),\qquad y(0)=\alpha,
\]
la récurrence devient
\[
\boxed{y_{k+1}=y_k-h\,t_k\,y_k^2}.
\]
La méthode reste identique : seule la fonction $f(t,y)$ change.

\section{Passage d'une équation d'ordre deux à un système d'ordre un}
Le schéma d'Euler s'applique directement à un système d'équations du premier ordre. Une équation d'ordre deux doit donc être transformée.

Pour le pendule proposé dans les sources,
\[
\ddot\theta(t)+\frac{g}{\ell}\sin\theta(t)=0,
\]
on pose
\[
y_0(t)=\theta(t),\qquad y_1(t)=\dot\theta(t).
\]
On obtient le système
\[
\begin{cases}
\dot y_0=y_1,\\
\dot y_1=-\dfrac{g}{\ell}\sin y_0.
\end{cases}
\]

\begin{center}
\TLSchemaEtatOrdreDeux
\end{center}

Le schéma d'Euler explicite devient
\[
\boxed{\begin{aligned}
y_{0,k+1}&=y_{0,k}+h y_{1,k},\\
y_{1,k+1}&=y_{1,k}-h\frac{g}{\ell}\sin(y_{0,k}).
\end{aligned}}
\]

\section{Implémentation en Python}
\begin{lstlisting}[style=TLPython]
def euler_explicit(f, t0, y0, h, n):
    """Retourne les listes des temps et des valeurs calculées."""
    t = [t0]
    y = [y0]
    for _ in range(n):
        y_next = y[-1] + h * f(t[-1], y[-1])
        t_next = t[-1] + h
        t.append(t_next)
        y.append(y_next)
    return t, y
\end{lstlisting}

Pour un système, $y_k$ est remplacé par un vecteur d'état et la relation
\[
\mathbf y_{k+1}=\mathbf y_k+h\,\mathbf f(t_k,\mathbf y_k)
\]
s'applique composante par composante.

\section{Choix du pas et validation}
Un pas plus petit augmente généralement la précision, mais accroît le nombre de calculs. La validation d'une simulation repose sur plusieurs vérifications :
\begin{itemize}
\item cohérence des unités et des conditions initiales ;
\item comparaison avec une solution exacte lorsqu'elle existe ;
\item répétition du calcul avec un pas plus petit ;
\item contrôle du comportement physique de la solution.
\end{itemize}

La méthode d'Euler explicite est simple et pédagogique, mais elle peut devenir imprécise ou instable lorsque le pas est trop grand.

\section{Synthèse de la démarche}
\begin{enumerate}
\item Mettre le modèle sous la forme $\dot{\mathbf y}=\mathbf f(t,\mathbf y)$.
\item Définir l'état initial $\mathbf y_0$, le pas $h$ et la durée de simulation.
\item Calculer successivement $\mathbf y_{k+1}=\mathbf y_k+h\mathbf f(t_k,\mathbf y_k)$.
\item Stocker les résultats et les représenter.
\item Vérifier leur cohérence et étudier l'influence du pas.
\end{enumerate}
